
% =========================================================================
%  Ejercicio 2b - Ensayo INDIVIDUAL sobre el articulo
%  Responsable: Carballido Camacateco Ana Lilia
%
%  NOTA: Este ensayo es INDIVIDUAL, uno por cada integrante del equipo.
% =========================================================================

\subsection*{Ensayo individual — Carballido Camacateco Ana Lilia}

\section*{Introducción}
El texto \cite{datavalue} presenta como tema principal la importancia que tienen los datos personales en la 
actualidad. Estos datos son generados todos los días y son de interés para cualquier sector empresarial ya que
tienen un gran impacto sobre las decisiones que toman con el fin de maximizar su beneficio económico.
Además, la autora nos hace conscientes sobre las multiples formas en las que generamos datos a través de nuestras 
interacciones con la tecnología. Dichas interacciones las hacemos a través del uso de redes sociales, aplicaciones,
servicios de streaming, entre otros. Por ejemplo, cuando usamos un servicio de streaming, este registra nuestras preferencias y gustos, lo que permite
a la empresa generar recomendaciones personalizadas y mejorar la experiencia del usuario. Sin embargo, también
implica que nuestros datos personales están siendo recopilados y utilizados con fines comerciales sin que nos
percatemos. \\
Actualmente, una gran parte de la población, somos usuaria de estos servicios de manera cotidiana a través de distintos
dispositivos como celulares, computadoras, relojes inteligentes, etc. Y, desde mi prespectiva, la mayoría de las
personas no son responsables con la información que comparten, ya que no se preocupan por leer los términos y condiciones de los servicios que utilizan.
Si bien las empresas son las principales responsables de proteger la información de sus usuarios, también es importante
que los usuarios sean conscientes de la información que comparten y cómo esta puede ser utilizada.

\vspace{0.3cm}

\section*{Desarrollo}
La autora tiene como objetivo dirigirnos a ser conscientes sobre la importancia que tienen nuestros datos personales, las
distintas formas en las que pueden ser utilizados y el valor que tienen para las empresas para así poder tomar decisiones informadas
y empoderarnos sobre nuestra información. Además, nos hace reflexionar sobre la importancia de ser responsables con la información que compartimos y cómo esta puede afectar nuestra privacidad y seguridad.
Este objetivo lo logra a través de distintos ejemplos en los que creemos que accedemos a un servicio gratuito, pero en realidad 
estamos cediendo nuestros datos personales a cambio de dicho servicio como sucede con el uso de redes sociales en las cuales se
hace un registro de nuestras interacciones, preferencias y gustos, lo que permite a las empresas generar publicidad dirigida y
mejorar la experiencia del usuario. Lo cual puede ser bueno en el aspecto de que encontramos contenido de nuestro interés, pero
también se han visto ejemplos en los cuales se ha hecho un mal uso de la información de los usuarios, como en el caso de Cambridge
Analytica, donde se utilizó la información de los usuarios de Facebook para influir en las elecciones presidenciales de Estados Unidos en 2016.\cite{rosenberg2018asi}\\
Por otra parte, hace incapié en el valor que nuestros datos personales tienen para las empresas en un sentido monetario. Nos cuenta
con cifras y ejemplos, cuanto dinero puede llegar a pagar una empresa por nuestros datos dependiendo del tipo de información que se 
pueda obtener de nosotros. Por ejemplo, una empresa puede pagar más por información sobre nuestros hábitos de consumo que por información
sobre nuestra ubicación geográfica. También influye la dificultad de obtener la información, ya que hay datos que los usuarios en promedio
prefieren no compartir como es el caso de la información sobre nuestra salud o finanzas personales.\\
La autora termina dando un ejemplo sobre una forma para aprovechar el internet de las cosas para la recolecciónde datos de manera
regulada y consensuada, en la cual los usuarios puedan decidir qué información quieren compartir y con quién, lo que permitiría a las empresas
obtener datos de manera más ética y responsable. \\
Como personas encargadas de la gestión de bases de datos, es importante que tengamos en cuenta la importancia de proteger la información 
de los usuarios y garantizar su privacidad. Además, debemos ser conscientes del valor que tienen los datos personales para las empresas
y cómo estos pueden ser utilizados para influir en nuestras decisiones y comportamientos. Por lo tanto, es fundamental que se implementen
medidas de seguridad y privacidad en la gestión de bases de datos, así como políticas claras sobre el uso y almacenamiento de la información de los usuarios.

\vspace{0.3cm}

\section*{Conclusión}
Como conclusión, podemos ver que la responsabilidad de proteger el uso e integridad denuestros datos personales recae
tanto en las empresas, que deben implementar estrategias de seguridad, privacidad y transparencia, como en los usuarios, que deben ser conscientes de la
información que comparten, cómo esta puede ser utilizada y el valor que tienen para las empresas.
Además, como profesionales en computación, también es una parte fundamental de nuestro trabajo poder garantizar la seguridad de los datos
que manejamos y proteger la información de los usuarios. Por lo tanto, es importante que nos mantengamos actualizados sobre las mejores prácticas 
en seguridad y privacidad de datos, así como en las regulaciones y leyes que rigen el uso de la información personal. \\






% Extension: entre 1 y 2 cuartillas por integrante.
%
% Estructura obligatoria (en parrafos, sin listas):
% - INTRODUCCION: presenta brevemente el tema, por que es importante, y
%   plantea tu postura o idea principal.
% - DESARROLLO: resume los argumentos del autor, relaciona con otros
%   conocimientos/lecturas/ejemplos, expone tus propios argumentos (a favor
%   o en contra) con razones, indica el objetivo que planteo el autor y como
%   se relaciona con la materia, y la tematica central y temas laterales y su
%   relacion con tu practica profesional.
% - CONCLUSION: reafirma tu postura, que aprendiste o la aportacion mas
%   valiosa, y si el tema abre preguntas o retos futuros.
% - FORMATO: parrafos completos, conectores logicos, ortografia y gramatica
%   cuidadas.

% Referencia del articulo (APA) - completar los datos:
% \cite{datavalue}